\documentclass[12pt,a4paper]{article}
\usepackage[margin=2.5cm]{geometry}
\usepackage[T1]{fontenc}
\usepackage[utf8]{inputenc}
\usepackage[brazil]{babel}
\usepackage{lmodern}
\usepackage{hyperref}
\usepackage{enumitem}
\usepackage{longtable}

\title{Sistema Instituto Alento \\ Documentacao Tecnica}
\author{Instituto Alento}
\date{\today}

\begin{document}
\maketitle
\tableofcontents
\newpage

\section{Visao geral tecnica}
\subsection{Stack}
\begin{itemize}[leftmargin=*]
\item Backend: Node.js + Express.
\item Frontend server-side: EJS + JS/CSS.
\item Banco: MongoDB com Mongoose.
\item Sessao: \texttt{express-session} com store Mongo.
\item Infra local/publica: Docker Compose + Nginx + Cloudflare Tunnel.
\end{itemize}

\subsection{Estrutura base}
\begin{itemize}[leftmargin=*]
\item \texttt{Controllers/}: controladores HTTP.
\item \texttt{Routes/}: rotas web e API.
\item \texttt{services/}: regras de negocio.
\item \texttt{schemas/}: modelos MongoDB.
\item \texttt{views/}: telas EJS.
\item \texttt{public/}: scripts e estilos frontend.
\end{itemize}

\section{Fluxo de autenticacao e autorizacao}
\begin{enumerate}
\item Usuario autentica no login.
\item Sessao segura e criada no servidor.
\item Middleware \texttt{requireAuth} protege rotas privadas.
\item Controle de acesso por permissao via \texttt{requirePermission}.
\item Rotas criticas de admin tambem usam \texttt{requireAdmin}.
\end{enumerate}

\section{Principais rotas funcionais}
\begin{longtable}{p{0.28\linewidth} p{0.68\linewidth}}
\textbf{Rota} & \textbf{Objetivo} \\
\hline
\texttt{/painel} & Dashboard principal da operacao. \\
\texttt{/painel/usuarios-online} & Lista usuarios online (admin). \\
\texttt{/agenda} & Calendario e gestao de agendamentos. \\
\texttt{/agenda/presencas} & Gestao de presencas e faltas. \\
\texttt{/familias} & Consulta e manutencao de familias/dependentes. \\
\texttt{/administracao} & Parametros operacionais (salas, justificativas, etc.). \\
\texttt{/api/administracao/*} & APIs de configuracao administrativa. \\
\end{longtable}

\section{Padrao de frontend}
\begin{itemize}[leftmargin=*]
\item JS modular por pagina (ex.: \texttt{public/js/administracao-actions.js}).
\item Validacao em tempo real para reduzir erro de envio.
\item Feedback inline por campo e feedback geral por formulario.
\item Renderizacao segura de dados dinamicos com sanitizacao/escape.
\end{itemize}

\section{Seguranca}
\begin{itemize}[leftmargin=*]
\item Protecao CSRF para operacoes sensiveis.
\item Headers de seguranca e telemetria de CSP.
\item Rate limit para reduzir abuso de requisicoes.
\item Testes focados em seguranca e autorizacao.
\end{itemize}

\section{Qualidade e testes}
\begin{itemize}[leftmargin=*]
\item Suite automatizada de seguranca (\texttt{npm test}).
\item Politica recomendada: nao publicar sem testes verdes.
\item Recomendado padrao de commit pequeno por tarefa.
\end{itemize}

\section{Deploy local com Docker}
\begin{enumerate}
\item Configurar \texttt{.env} a partir de \texttt{.env.example}.
\item Executar \texttt{docker compose up -d --build}.
\item Validar app em \texttt{/login}.
\end{enumerate}

\section{Guia de manutencao}
\subsection{Quando adicionar um modulo novo}
\begin{itemize}[leftmargin=*]
\item Criar rotas novas em \texttt{Routes/}.
\item Criar controller dedicado em \texttt{Controllers/}.
\item Concentrar regra de negocio em \texttt{services/}.
\item Adicionar tela EJS e JS da pagina.
\item Incluir permissao no controle de acesso.
\item Cobrir cenarios criticos com teste.
\end{itemize}

\subsection{Quando alterar modulo existente}
\begin{itemize}[leftmargin=*]
\item Evitar quebrar contratos de API ja usados.
\item Revisar impacto em perfis de acesso.
\item Validar navegacao em desktop e mobile.
\end{itemize}

\end{document}
